\documentclass[spanish,letterpaper,oneside]{article}

\def\documenttitle {Elementos del delito de robo}
\def\documentsubtitle {Cuadro comparativo y aplicacion al caso}
\def\documentsubject {Actividad 2 - Derecho penal especial mexicano}
\def\documentauthor {Martin Jonathan de la Cruz}
\def\coursename {Derecho penal especial mexicano}
\def\coursecode {020}
\def\universityname {Universidad Abierta y a Distancia de Mexico}
\def\universityfaculty {Licenciatura en Derecho}
\def\universitydepartment {Derecho penal especial mexicano}
\def\universitydepartmentimage {departamentos/UnADM}
\def\universitydepartmentimagecfg {height=1.57cm}
\def\universitylocation {Roma Norte, Ciudad de Mexico}
\def\authortable {
  \begin{tabular}{ll}
    \textbf{Alumno:} & Martin Jonathan de la Cruz \\
    Matricula: & ES2611202040 \\
    Figura docente: & Veronica Ortiz Lopez \\
    Semestre/Bloque: & 2 / 2 \\
    Semana: & 2 \\
    \multicolumn{2}{l}{Fecha de realizacion: \today} \\
    \multicolumn{2}{l}{\universitylocation}
  \end{tabular}
}

\input{template}
\usepackage{helvet}
\renewcommand{\familydefault}{\sfdefault}
\usepackage{array}
\usepackage{booktabs}
\usepackage{longtable}
\usepackage{pdflscape}
\usepackage{calc}
\usepackage{tikz}
\usetikzlibrary{arrows.meta,positioning,calc,babel}
\setcitestyle{authoryear,open={(},close={)}}
\providecommand{\tightlist}{\setlength{\itemsep}{0pt}\setlength{\parskip}{0pt}}
\providecommand{\pandocbounded}[1]{#1}
\providecommand{\passthrough}[1]{#1}
\setlength{\emergencystretch}{3em}

\def\coverwatermarkenabled {true}
\def\coverwatermarkimage {img/departamentos/UnADM.pdf}
\def\coverwatermarkopacity {0.16}
\def\coverwatermarkwidth {0.72\paperwidth}
\def\coverwatermarkxshift {0cm}
\def\coverwatermarkyshift {-0.35cm}
\newcommand{\insertcoverwatermark}{%
  \ifthenelse{\equal{\coverwatermarkenabled}{true}}{%
    \AddToShipoutPictureBG*{%
      \begin{tikzpicture}[remember picture,overlay]
        \node[opacity=\coverwatermarkopacity,inner sep=0pt,
          xshift=\coverwatermarkxshift,yshift=\coverwatermarkyshift]
          at (current page.center) {%
            \includegraphics[width=\coverwatermarkwidth]{\coverwatermarkimage}%
          };
      \end{tikzpicture}%
    }%
  }{}%
}

\begin{document}
\insertcoverwatermark
\templatePortrait
\templatePagecfg
\onehalfspacing
\templateIndex
\templateFinalcfg

\section{Introduccion}\label{introduccion}

La parte especial del derecho penal identifica conductas que afectan
bienes juridicos concretos. Comparar el robo exige distinguir el objeto
material del interes protegido y reconocer cada elemento normativo; la
semejanza entre dos disposiciones no implica que ambas se apliquen
simultaneamente al mismo hecho. La comparacion del supuesto del
articulo 367 del Codigo Penal Federal (CPF) con el del articulo 220 del
Codigo Penal para el Distrito Federal (CPDF), utilizando las copias
proporcionadas en el aula y el material de apoyo S2 permite distinguir
la formulacion normativa y su aplicacion al caso del telefono celular
\citep{unadmPenalS2}.

\section{Elementos del robo y aplicacion juridica}
\subsection{Lectura del tipo penal}
El cuadro diferencia norma, sujetos, conducta y bien juridico. El
patrimonio es el interes protegido; el telefono constituye el objeto
material. La comparacion se limita al supuesto basico, sin fijar penas
ni decidir la competencia del tribunal. Cada elemento debe relacionarse
con los hechos, no deducirse unicamente de la ausencia del propietario.

\begin{landscape}
\thispagestyle{empty}
\subsection{Cuadro comparativo}\label{cuadro-comparativo}
\begingroup
\setstretch{1.05}
\footnotesize
\setlength{\tabcolsep}{5pt}
\renewcommand{\arraystretch}{1.15}
\noindent\begin{tabular}{@{}
  >{\raggedright\arraybackslash}p{0.16\linewidth}
  >{\raggedright\arraybackslash}p{0.40\linewidth}
  >{\raggedright\arraybackslash}p{\dimexpr0.44\linewidth-4\tabcolsep\relax}@{}}
\toprule
\textbf{Elemento} & \textbf{Codigo Penal Federal} &
\textbf{Codigo Penal para el Distrito Federal} \\
\midrule
Articulo aplicable & 367; 369 solo para precisar consumacion. & 220,
primer parrafo; 222 solo para contrastar uso y dominio. \\
Descripcion legal & ``Comete el delito de robo: el que se apodera de una
cosa ajena mueble, sin derecho y sin consentimiento de la persona que
puede disponer de ella con arreglo a la ley.'' & ``Al que con animo de
dominio y sin consentimiento de quien legalmente pueda otorgarlo, se
apodere de una cosa mueble ajena, se le impondran:''. Se transcribe el
supuesto del primer parrafo; se omiten las penas porque no son objeto de
la comparacion. \\
Bien juridico tutelado & El patrimonio, afectado por la privacion del
control juridicamente protegido sobre la cosa; no confundirlo con el
telefono como objeto fisico. & El patrimonio, protegido frente al
apoderamiento no consentido de una cosa ajena. \\
Sujeto activo & En el supuesto basico, cualquier persona que realiza el
apoderamiento; el articulo no exige una calidad especial. & Cualquier
persona que realiza el apoderamiento con animo de dominio; no exige
calidad especial para este supuesto. \\
Sujeto pasivo & Titular del interes patrimonial afectado; en el caso
planteado, el trabajador propietario del telefono. La persona facultada
para consentir puede ser relevante por su facultad juridica de
disposicion. & Titular del interes patrimonial lesionado; en el caso, el
trabajador propietario. El consentimiento debe provenir de quien pueda
otorgarlo legalmente. \\
Conducta delictiva & Apoderarse: colocar la cosa bajo el propio control
sin derecho y sin autorizacion juridicamente valida. & Apoderarse: tomar
control sobre la cosa sin consentimiento valido y con proposito de
comportarse como dueno. \\
Elementos del tipo & Apoderamiento; cosa mueble; ajenidad; ausencia de
derecho; falta de consentimiento de persona facultada. Deben
acreditarse, no presumirse solo por la ausencia del propietario. &
Apoderamiento; cosa mueble; ajenidad; falta de consentimiento de persona
facultada; animo de dominio expresamente mencionado. \\
\bottomrule
\end{tabular}
\captionof{table}{Comparacion del supuesto basico de robo: articulos 367 CPF y 220 CPDF.}

Fuentes normativas: \citep[art.~367]{cpfAula2026} y
\citep[art.~220]{cpdfAulaS2}. Explicaciones contrastadas con
\citep{unadmPenalS2}. La tabla no determina vigencia de penas ni
jurisdiccion competente.
\endgroup
\end{landscape}

\subsection{Coincidencias y diferencias}\label{coincidencias-y-diferencias}

\textbf{Coincidencia 1:} ambas disposiciones exigen apoderamiento de una
cosa mueble ajena. No basta tener contacto con el objeto: debe existir
toma de control sobre algo que no pertenece al sujeto
\citep[art.~367]{cpfAula2026}\citep[art.~220]{cpdfAulaS2}.

\textbf{Coincidencia 2:} ambas requieren ausencia de consentimiento de
quien pueda disponer o consentir conforme al derecho. La ausencia fisica
del propietario no equivale por si misma a falta de permiso; el supuesto
la acredita expresamente al negar toda autorizacion
\citep[art.~367]{cpfAula2026}\citep[art.~220]{cpdfAulaS2}.

\textbf{Diferencia 1:} el articulo 220 expresa el ``animo de dominio'';
el 367 no utiliza esas palabras. Es una diferencia de formulacion, no
una prueba de que el robo federal carezca de componente subjetivo ni de
que sea una conducta accidental.

\textbf{Diferencia 2:} el articulo 367 incluye literalmente ``sin
derecho'' ademas de falta de consentimiento. El primer parrafo del 220
no reproduce esa expresion como segmento separado. No cabe deducir que
la legislacion local castigue el ejercicio licito de un derecho; la
comparacion identifica el tenor normativo, no elimina las reglas
generales de antijuridicidad. Por tanto, los elementos centrales
examinados son ampliamente coincidentes y las diferencias descritas no
autorizan conclusiones absolutas sobre toda la regulacion.

\subsection{Aplicacion al caso del telefono
celular}\label{aplicacion-al-caso-del-telefono-celular}

\textbf{a) Conducta juridicamente relevante.} Tomar el telefono de otro
trabajador, guardarlo entre las pertenencias propias y retirarse para
conservarlo. Estos actos describen un apoderamiento, no la mera
observacion del celular; el dato de no contar con autorizacion excluye
un prestamo consentido en los hechos planteados.

\textbf{b) Sujetos.} Sujeto activo: la persona que toma el dispositivo.
Sujeto pasivo: el trabajador al que pertenece. La oficina es el lugar,
no el sujeto pasivo; no hay base para afirmar que el empleador sea
propietario del telefono.

\textbf{c) Bien juridico.} El patrimonio del trabajador. El objeto
material es el telefono. Esta distincion evita confundir el interes
protegido por el derecho con el objeto sobre el cual recae la accion
\citep{unadmPenalS2}.

\textbf{d) Elementos identificables.} Apoderamiento: tomar y llevar el
dispositivo; movilidad: un celular puede trasladarse; ajenidad:
pertenece al trabajador; falta de consentimiento: no se solicito ni
obtuvo autorizacion; falta de derecho: el supuesto no atribuye facultad
juridica alguna al agente. La intencion de conservarlo permite
identificar el animo de dominio exigido expresamente por el articulo
220. El articulo 369 CPF situa la consumacion en tener la cosa en poder
del agente, por lo que no es indispensable venderla para configurar ese
aspecto.

\textbf{e) Correspondencia con ambas legislaciones.} En un ejercicio de
subsuncion de los hechos dados, se reconocen los elementos del articulo
367 CPF y del primer parrafo del 220 CPDF: apoderamiento de cosa mueble
ajena, ausencia de derecho o autorizacion y, en el texto local, dominio.
La intencion de conservarlo distingue el caso del uso temporal al que
alude el articulo 222 CPDF. Esta conclusion es provisional y academica:
no sustituye la prueba procesal, no decide culpabilidad ni determina
competencia federal/local. No se conoce valor del telefono,
circunstancias completas ni causa de justificacion que permita fijar una
pena.

\clearpage
\section{Conclusiones}\label{conclusion}

Comparar tipos penales requiere separar su texto de las inferencias
sobre el caso. El nucleo compartido es el apoderamiento no consentido de
una cosa mueble ajena; las expresiones ``sin derecho'' y ``animo de
dominio'' obligan a justificar el analisis, no a inventar una oposicion
total entre codigos. La aplicacion al celular muestra que sujetos,
objeto y bien juridico cumplen funciones diferentes. Un razonamiento
riguroso identifica elementos reconocibles y tambien limites probatorios
y competenciales.
Considero que la precision importa mas que afirmar una diferencia
artificial entre normas. Desde mi formacion en Derecho, prefiero
justificar cada elemento y reconocer los datos que faltan antes de
concluir sobre responsabilidad o pena. Esa cautela permite formular
argumentos revisables sin sustituir la prueba por una intuicion.\footnote{
Se utilizo GitHub Copilot para apoyar la organizacion, la redaccion y
la maquetacion. La asistencia no sustituye el analisis propio ni la
revision personal de las fuentes y los argumentos antes de entregar.}

\clearpage
\begin{thebibliography}{3}
\bibitem[Asamblea Legislativa del Distrito Federal(s.\,f.)]{cpdfAulaS2}
Asamblea Legislativa del Distrito Federal. (s.\,f.). \emph{Codigo Penal
para el Distrito Federal} (arts. 220 y 222, copia proporcionada en el
aula). Instituto de Investigaciones Legislativas.
\url{https://data.consejeria.cdmx.gob.mx/index.php/leyes/codigos/34-codigo-penal-para-el-distrito-federal}

\bibitem[Camara de Diputados del H. Congreso de la Union(2026)]{cpfAula2026}
Camara de Diputados del H. Congreso de la Union. (2026). \emph{Codigo
Penal Federal} (ultima reforma indicada en la copia: DOF 13-03-2026;
arts. 367 y 369). \url{https://www.diputados.gob.mx/LeyesBiblio/pdf/CPF.pdf}

\bibitem[Universidad Abierta y a Distancia de Mexico(2026)]{unadmPenalS2}
Universidad Abierta y a Distancia de Mexico. (2026). \emph{Material de
apoyo S2: Derecho penal especial mexicano} {[}Material de aula, recurso
215278{]}.
\url{https://aulavirtual.unadmexico.mx/mod/resource/view.php?id=215278}
\end{thebibliography}

\end{document}